\documentclass[12pt]{article}
\usepackage[utf8]{inputenc}
\usepackage[T1]{fontenc}
\usepackage{hyperref}
\title{Pol Language Resources for K-12 Curriculum Development}
\author{Compiled by Claude Sonnet 5.5 (claude-sonnet-5-5) for Andr\'as Kornai}
\date{}
\begin{document}
\maketitle

\section{Introduction}
Pol (ISO 639-3 pmm; Glottocode pomo1271) is a Bantu language of Cameroon, with a smaller community reported across the border. Sources disagree on speaker numbers (Wikipedia cites about 44,000 for Pol and Pomo speakers; Joshua Project gives a larger ethnic population). Very little educational material turned up in search; this dossier lists only reference sources.

\section{Primary Educational Resources}
No primers, textbooks or school programmes for Pol were found in search results.

\section{Linguistic and Reference Materials}
\begin{itemize}
\item Koeppe, \emph{Grammaire descriptive du polri} (descriptive grammar covering phonology, morphology and syntax; bookseller page): \url{https://www.koeppe.de/titel_print_grammaire-descriptive-du-polri}
\item Wikipedia, Pol language: \url{https://en.wikipedia.org/wiki/Pol_language}
\item Glottolog, Pol (pomo1271): \url{https://glottolog.org/resource/languoid/id/pomo1271.bigmap.html}
\end{itemize}

\section{Religious and Community Resources}
\begin{itemize}
\item Joshua Project, Pol in Cameroon (people-group profile): \url{https://joshuaproject.net/index.php/people_groups/14429/CM}
\end{itemize}
No scripture translations were found in the search results.

\section{Digital and Audio Resources}
No audio, video, corpus or software resources for Pol were found.

\section{Curriculum Development Context}
The only substantial documentation located is the descriptive grammar above. A native-language curriculum would need new primers and orthographic work, ideally with the community; nothing found indicates existing school use.
\end{document}